% !TeX root = ../../book.tex
\chapter*{前言}
\addcontentsline{toc}{chapter}{前言}
\markboth{前言}{前言}

你好，感谢你花时间阅读\textit{《纯数学的无穷递降》}一书的快速导览！本书的最新版本可以从下面网站免费下载：
\begin{center} \vspace{-10pt}
\url{\bookurl}
\end{center} \vspace{-10pt}
这个网站还包含本书不同版本间的更新信息，历史版本归档，以及使用 \LaTeX{} 的一些资源（参见 \Cref{apxLaTeX}）。

\subsection*{关于本书}

在典型的微积分课上，学生会学习链式法则，然后运用该法则去解决一些给定的`链式法则问题'，比如计算 $\sin(1+x^2)$ 对 $x$ 的导数，或者解决涉及变化率的应用题。期望学生正确应用链式法则得出正确答案，并给出完整的步骤以确信结论正确。从这个意义上说，学生是数学的\textit{消费者}。他们被赋予链式法则这个必须无条件接受的工具，然后使用这个工具来解决范围狭窄的问题。

本书的目标是帮助读者从数学的\textit{消费者}转变成数学的\textit{生产者}。这就是`纯'数学的意思。数学的消费者可能会学习链式法则并用它来计算导数，而数学的生产者可能会从导数的严格定义中推导出链式法则，然后在更一般的情况下（例如多元分析）证明链式法则的更抽象版本。

数学的消费者只需要说出他们是如何使用工具找到答案的。而数学的生产者必须做更多事情：他们必须能够追踪定义和假设，以新颖有趣的方式拼凑事实，并对数学概念给出自己的定义。更重要的是，一旦这样做了，他们必须以其他人可以理解的方式交流他们的发现，并且必须说服其他人他们所做的是正确的、适当的和有价值的。

正是这种从数学消费到生产的转变指导了我用于设计和编写本书的原则。尤其是：

\begin{itemize}
\item \textbf{交流  } 首先，本书旨在帮助读者获得数学素养并用数学的方式表达自己。这发生在许多层次上，例如，考虑以下表达式：。
\[ \forall x \in \mathbb{R},\, [\neg(x = 0) \Rightarrow (\exists y \in \mathbb{R},\, y^2 < x^2)] \]
学完本书，你将能够直观地说出符号 $\forall$, $\in$, $\mathbb{R}$, $\neg$, $\Rightarrow$ 和 $\exists$ 的含义以及它们的精确定义。你也能够从整体上解释这个表达的含义，用简单的术语向他人解释它的意思而\textit{不是}使用一堆混乱的符号，你能够证明它是正确的，并用清晰简明的方式与他人交流你的证明。

为了实现此目标，本书主要章节都在构建这类工具，而 \Cref{apxWriting} 会更加关注书写层面。

\item \textbf{探究  } 研究表明，人们在自己发现事物时会学到更多。如果我把它发挥到极致，那这本书将是一片空白；然而，我的确认为将探究式学习的各个方面纳入本书很重要。

这一原则体现在本书中散布的习题，其中许多习题只是要求你证明给出的结果。许多读者会觉得这令人沮丧，但这是有充分理由的：这些习题作为检查点，确保你对内容的理解足以继续前行。那种挫败感就是学习的感觉 --- 请拥抱它！

\item \textbf{策略  } 数学证明很像一个谜题。在证明的任意给定阶段，你都会有若干可用的定义、假设和结果，你必须使用你可以支配的逻辑规则将它们拼凑在一起。贯穿全书，尤其是前面的章节，当有用的证明策略出现时，我不遗余力地将它们强调出来。

\item \textbf{内容  } 如果你对要证明的结果没有任何概念，那么学习数学就没有多大意义。考虑到这一点，\Cref{ptTopics} 包含几个章节专门介绍纯数学中的一些主题领域，例如数论、组合数学、分析和概率论。

\item \textbf{\LaTeX{}  } 数学排版的\textit{事实}标准是 \LaTeX{}。我认为数学家在有指导的情况下尽早学习 \LaTeX{} 这一点很重要，因此我在 \Cref{apxLaTeX} 中编写了一个简短的教程，并且全书所有新符号在定义时都给出了 \LaTeX{} 代码。

\end{itemize}



\subsection*{浏览本书}

这本书不需要，而且，重点强调一下，也\textit{不应该}，从前到后阅读。选择内容呈现顺序是为了使后面出现的内容仅依赖于前面出现的内容，但是按照内容的呈现顺序进行学习可能是一种相当枯燥的体验。

大多数纯数学导论课程至少涵盖符号逻辑、集合、函数和关系。这是 \Cref{ptCoreConcepts} 的内容。此类课程通常涵盖纯数学的附加主题，具体主题取决于课程为学生准备的内容。考虑到这一点，\Cref{ptTopics} 介绍了纯数学的一系列领域，包括数论、组合学、集合论、实分析、概率论和阶论。

在继续讨论 \Cref{ptTopics} 中的主题之前，没有必要涵盖所有 \Cref{ptCoreConcepts} 的内容。事实上，穿插 \Cref{ptTopics} 中的内容可能是激发 \Cref{ptCoreConcepts} 中出现的许多抽象概念的有用方法。

下表显示了章节之间的依赖关系。除非另有说明，否则同一章节中的前几节均被视为`基本'先决条件。

\begin{center}
\begin{tabular}{c|c|ccc}
\textbf{部分} & \textbf{章节} & \textbf{必修} & \textbf{建议} & \textbf{重点} \\ \hline
\textbf{\ref{ptCoreConcepts}} & \ref{secPropositionalLogic} & \ref{chGettingStarted} &  &  \\
&\ref{secSets} & \ref{secLogicalEquivalence} &  &  \\
&\ref{secFunctions} & \ref{secSetOperations} &  &  \\
&\ref{secPeanosAxioms} & \ref{secLogicalEquivalence} & \ref{secFunctions} & \ref{secInjectionsSurjections} \\
&\ref{secRelations} & \ref{secSets} & \ref{secFunctions} & \ref{secInjectionsSurjections}, \ref{secWeakInduction} \\
&\ref{secFiniteSets} & \ref{secInjectionsSurjections}, \ref{secStrongInduction} & \ref{secEquivalenceRelationsPartitions} &  \\ \hline
\textbf{\ref{ptTopics}} & \ref{secDivision} & \ref{secLogicalEquivalence} & \ref{secSets}, \ref{secStrongInduction} & \ref{secFunctions} \\
&\ref{secModularArithmetic} &  & \ref{secEquivalenceRelationsPartitions} &  \\
&\ref{secCountingPrinciples} & \ref{secFiniteSets} &  \\
&\ref{secInequalitiesMeans} & \ref{secPeanosAxioms}, \ref{secSets} &  & \ref{secEquivalenceRelationsPartitions} \\
&\ref{secCompletenessConvergence} & \ref{secFunctions} & \ref{secInequalitiesMeans} &  \\
&\ref{secSeriesSums} & \ref{secFunctions} & \ref{secInequalitiesMeans} & \ref{secModularArithmetic}, \ref{secCountableUncountableSets} \\
&\ref{secCardinality} & \ref{secCountableUncountableSets} &  &  \\
&\ref{secCardinalArithmetic} & \ref{secCountingPrinciples} &  &  \\
&\ref{secDiscreteProbabilitySpaces} & \ref{secCountingPrinciples} & \ref{secCountableUncountableSets}, \ref{secSeriesSums} &  \\
&\ref{secOrdersLattices} & \ref{secEquivalenceRelationsPartitions} &  &  \\
&\ref{secStructuralInduction} & \ref{secStrongInduction}, \ref{secInjectionsSurjections} & \ref{secCountableUncountableSets} & \ref{secOrdersLattices}
\end{tabular}
\end{center}

先决条件是累积的。例如，为了学习 \Cref{secCardinalArithmetic}，你应该首先学习 \Crefrange{chGettingStarted}{chMathematicalInduction} 和 \Cref{secFiniteSets,secCountingPrinciples,secCountableUncountableSets,secCardinality}。

\subsection*{数字、颜色和符号的含义}

广义上讲，本书中的内容被分解到五类项目中：定义、结论、备注、示例和练习。在 \Cref{apxWriting} 中，我们也有证明摘录。为了提高导航性，这些类别通过名称、颜色和符号进行区分，如下表所示。

\begin{center}
\begin{tabular}{lcl}
\textbf{类型} & \textbf{符号} & \textbf{颜色} \\ \hline
定义 & \defintrosymbol & {\fontfamily{bch}\color{defcol} \textbf{红色}} \\
结论 & \thmintrosymbol & {\fontfamily{bch}\color{thmcol} \textbf{蓝色}} \\
备注 & \tipintrosymbol & {\fontfamily{bch}\color{tipcol} \textbf{紫色}} \\
\end{tabular}
\hspace{20pt}
\begin{tabular}{lcl}
\textbf{类型} & \textbf{符号} & \textbf{颜色} \\ \hline
示例 & \exintrosymbol & {\fontfamily{bch}\color{excol} \textbf{蓝绿色}} \\
练习 & \printrosymbol & {\fontfamily{bch}\color{prcol} \textbf{金色}} \\
证明摘录 & \quoteintrosymbol & {\fontfamily{bch}\color{excol} \textbf{蓝绿色}}
\end{tabular}
\end{center}

这些项目根据它们所在的章节逐一呈现出来 --- 例如，\Cref{thmUniquenessofLimits} 在 \Cref{secCompletenessConvergence} 中。定义和定理（重要结果）出现在 \fbox{方框} 中。

你还会遇到符号 $\square$ 和 \nonproofqedsymbol{} 它们的含义如下：

\begin{itemize}
\item[$\square$] \textbf{证明结束  } 这是数学文档中的标准，通过画一个小正方形或写下`\textit{Q.E.D.}'（后者代表 \textit{quod erat demonstrandum}，在拉丁语中表示 \textit{这是要展示的}。）表示证明结束。
\item[\nonproofqedsymbol] \textbf{条目结束  } 这\textit{不}是标准用法，将其包含在内只是为了帮助你识别项目结束，接下来是本书的正文内容。
\end{itemize}

一些小节标有符号 \optmarksymbol{}。这表明可以跳过该小节中的内容而不会产生严重后果。

\subsection*{许可}

本书根据 Creative Commons Attribution-ShareAlike 4.0 (\textsc{cc by-sa 4.0}) 许可证获得许可。这意味着你可以分享本书，但前提是你要注明作者的姓名，并且本书的任何副本或衍生品都是在同一许可下发布的。

许可的完整全文附在本书末尾或者通过下方链接阅读：

\begin{center}
\url{http://creativecommons.org/licenses/by-sa/4.0/}
\end{center}

\subsection*{评论和修正}

任何反馈，无论是来自学生、助教、教师还是任何其他读者，我们都不胜感激。尤其是印刷错误的修正、概念描述或定理证明的建议替代方法，以及对新内容的需求（例如，如果你知道一个说明某个概念的绝佳例子，或者某个你希望本书涵盖的相关概念）。

此类反馈可以通过电子邮件分别发送给作者\ifadapted{和译者} (\url{\authoremail}\ifadapted{ 和 \url{\adapteremail}}).